% ============================================================
%  MTH636 -- Game Theory | Lecture 11
%  Computing Mixed Strategies (Graphical Method)
% ============================================================
\documentclass[12pt, a4paper]{article}

% ============================================================
%  MTH636 -- Game Theory
%  Notes Preamble — Color-Coded, TikZ-Enhanced
% ============================================================

% --- Core packages ---
\usepackage[T1]{fontenc}
\usepackage[utf8]{inputenc}
\usepackage{lmodern}
\usepackage[margin=2.5cm, top=2.8cm, bottom=2.8cm]{geometry}
\usepackage{microtype}

% --- Math ---
\usepackage{amsmath, amssymb, amsthm, mathtools}
\usepackage{bm}        % bold math
\usepackage{cancel}    % strike-through in math

% --- Color & Boxes ---
\usepackage[dvipsnames, svgnames, x11names, table]{xcolor}
\usepackage[most]{tcolorbox}
\tcbuselibrary{breakable, skins, theorems}

% --- TikZ & Diagrams ---
\usepackage{tikz}
\usepackage{tikz-qtree}
\usetikzlibrary{arrows.meta, positioning, shapes.geometric, fit, backgrounds,
                decorations.pathreplacing, calc, matrix, trees}
\usepackage{pgfplots}
\pgfplotsset{compat=1.18}

% --- Tables ---
\usepackage{booktabs}
\usepackage{multirow}
\usepackage{array}
\usepackage{tabularx}
\usepackage{makecell}

% --- Typography & Layout ---
\usepackage{enumitem}
\usepackage{parskip}
\usepackage{fancyhdr}
\usepackage{titlesec}
\usepackage{titletoc}
\usepackage{hyperref}
\usepackage{graphicx}
\usepackage{float}
\usepackage{caption}

% --- Color Palette ---
\definecolor{defblue}{HTML}{1A5276}      % Definition header
\definecolor{defbluebg}{HTML}{D6EAF8}    % Definition background
\definecolor{thmgreen}{HTML}{1E8449}     % Theorem header
\definecolor{thmbg}{HTML}{D5F5E3}        % Theorem background
\definecolor{exorange}{HTML}{B7770D}     % Example header
\definecolor{exbg}{HTML}{FEF9E7}         % Example background
\definecolor{warnred}{HTML}{922B21}      % Warning/Key Insight header
\definecolor{warnbg}{HTML}{FADBD8}       % Warning background
\definecolor{noticegray}{HTML}{2C3E50}   % Note header
\definecolor{noticebg}{HTML}{EBF5FB}     % Note background
\definecolor{taketawayyellow}{HTML}{7D6608} % Key Takeaway header
\definecolor{takeawaybg}{HTML}{FDFEFE}   % Takeaway background
\definecolor{accentpurple}{HTML}{6C3483} % Accent / proof
\definecolor{lightgray}{HTML}{F2F3F4}    % shading

% ============================================================
%  TCOLORBOX ENVIRONMENTS
% ============================================================

% -- Definition (Blue) --
\newtcolorbox{defbox}[2][]{
  enhanced, breakable,
  colback=defbluebg, colframe=defblue, coltitle=white,
  fonttitle=\bfseries\sffamily,
  title={Definition: #2},
  attach boxed title to top left={yshift=-2mm, xshift=4mm},
  boxed title style={colback=defblue, rounded corners},
  arc=4pt, boxrule=1.2pt,
  left=6pt, right=6pt, top=8pt, bottom=6pt,
  #1
}

% -- Theorem / Proposition (Green) --
\newtcolorbox{thmbox}[2][]{
  enhanced, breakable,
  colback=thmbg, colframe=thmgreen, coltitle=white,
  fonttitle=\bfseries\sffamily,
  title={Theorem: #2},
  attach boxed title to top left={yshift=-2mm, xshift=4mm},
  boxed title style={colback=thmgreen, rounded corners},
  arc=4pt, boxrule=1.2pt,
  left=6pt, right=6pt, top=8pt, bottom=6pt,
  #1
}

% -- Proposition (Green variant) --
\newtcolorbox{propbox}[2][]{
  enhanced, breakable,
  colback=thmbg, colframe=thmgreen!70!black, coltitle=white,
  fonttitle=\bfseries\sffamily,
  title={Proposition: #2},
  attach boxed title to top left={yshift=-2mm, xshift=4mm},
  boxed title style={colback=thmgreen!70!black, rounded corners},
  arc=4pt, boxrule=1.2pt,
  left=6pt, right=6pt, top=8pt, bottom=6pt,
  #1
}

% -- Example (Orange/Yellow) --
\newtcolorbox{exbox}[2][]{
  enhanced, breakable,
  colback=exbg, colframe=exorange, coltitle=white,
  fonttitle=\bfseries\sffamily,
  title={Example: #2},
  attach boxed title to top left={yshift=-2mm, xshift=4mm},
  boxed title style={colback=exorange, rounded corners},
  arc=4pt, boxrule=1.2pt,
  left=6pt, right=6pt, top=8pt, bottom=6pt,
  #1
}

% -- Key Insight / Warning (Red) --
\newtcolorbox{warnbox}[2][]{
  enhanced, breakable,
  colback=warnbg, colframe=warnred, coltitle=white,
  fonttitle=\bfseries\sffamily,
  title={Key Insight: #2},
  attach boxed title to top left={yshift=-2mm, xshift=4mm},
  boxed title style={colback=warnred, rounded corners},
  arc=4pt, boxrule=1.2pt,
  left=6pt, right=6pt, top=8pt, bottom=6pt,
  #1
}

% -- Note (Gray-Blue) --
\newtcolorbox{notebox}[1][]{
  enhanced, breakable,
  colback=noticebg, colframe=noticegray!60, coltitle=white,
  fonttitle=\bfseries\sffamily,
  title={Note},
  attach boxed title to top left={yshift=-2mm, xshift=4mm},
  boxed title style={colback=noticegray!60, rounded corners},
  arc=4pt, boxrule=0.8pt,
  left=6pt, right=6pt, top=8pt, bottom=6pt,
  #1
}

% -- Key Takeaway --
\newtcolorbox{takebox}[1][]{
  enhanced, breakable,
  colback=takeawaybg, colframe=taketawayyellow, coltitle=white,
  fonttitle=\bfseries\sffamily,
  title={Key Takeaway},
  attach boxed title to top left={yshift=-2mm, xshift=4mm},
  boxed title style={colback=taketawayyellow, rounded corners},
  arc=4pt, boxrule=1.2pt,
  left=6pt, right=6pt, top=8pt, bottom=6pt,
  #1
}

% -- Intuition --
\newtcolorbox{intubox}[1][]{
  enhanced, breakable,
  colback=accentpurple!8, colframe=accentpurple!60, coltitle=white,
  fonttitle=\bfseries\sffamily,
  title={Intuition},
  attach boxed title to top left={yshift=-2mm, xshift=4mm},
  boxed title style={colback=accentpurple!80, rounded corners},
  arc=4pt, boxrule=0.8pt,
  left=6pt, right=6pt, top=8pt, bottom=6pt,
  #1
}

% ============================================================
%  PAYOFF MATRIX HELPER
%  Usage: inside tabular, use \payoff{a}{b} for (a,b)
% ============================================================
\newcommand{\payoff}[2]{$(#1,\;#2)$}
\newcommand{\nash}[2]{\underline{$#1$},\underline{$#2$}}  % underline NE payoffs
\newcommand{\Nset}{\mathbb{N}}
\newcommand{\Rset}{\mathbb{R}}
\newcommand{\Sset}{\mathcal{S}}
\newcommand{\Ui}{u_i}

% ============================================================
%  SECTION FORMATTING
% ============================================================
\titleformat{\section}
  {\large\bfseries\sffamily\color{defblue}}
  {\thesection}{1em}{}[\titlerule]

\titleformat{\subsection}
  {\normalsize\bfseries\sffamily\color{thmgreen!80!black}}
  {\thesubsection}{1em}{}

\titleformat{\subsubsection}
  {\normalsize\itshape\sffamily\color{exorange!80!black}}
  {\thesubsubsection}{1em}{}

% ============================================================
%  HEADER / FOOTER
% ============================================================
\pagestyle{fancy}
\fancyhf{}
\fancyhead[L]{\sffamily\small\color{defblue} MTH636 -- Game Theory}
\fancyhead[R]{\sffamily\small\color{defblue} \leftmark}
\fancyfoot[C]{\sffamily\small\thepage}
\renewcommand{\headrulewidth}{0.4pt}

% ============================================================
%  HYPERREF SETUP
% ============================================================
\hypersetup{
  colorlinks=true,
  linkcolor=defblue,
  urlcolor=thmgreen,
  citecolor=accentpurple,
  bookmarksnumbered=true,
  pdftitle={MTH636 Game Theory Notes},
  pdfauthor={MTH636 Student Notes}
}

% ============================================================
%  MISC HELPERS
% ============================================================
\newcommand{\BR}{\text{BR}}        % Best Response
\newcommand{\NE}{\text{NE}}        % Nash Equilibrium
\newcommand{\SPNE}{\text{SPNE}}    % Subgame Perfect NE
\newcommand{\Dom}{\text{dom}}      % Dominance
\newcommand{\mixs}[1]{\sigma_{#1}} % Mixed strategy
\newcommand{\pures}[1]{s_{#1}}     % Pure strategy


% ============================================================
\begin{document}

% Title
\begin{center}
  {\LARGE\bfseries\sffamily\color{defblue} MTH636 -- Game Theory}\\[6pt]
  {\Large\sffamily Lecture 11: Computing Max-Min Strategies graphically}\\[4pt]
\end{center}
\vspace{0.5cm}
\hrule height 1.5pt
\vspace{0.3cm}

% ============================================================
\section{Max-Min Strategies in Mixed Games}

For a finite two-player zero-sum game, we know the value $V$ exists in mixed strategies:
\[
  V = \max_{\sigma_1 \in \Sigma_1} \min_{s_2 \in S_2} U(\sigma_1, s_2) = \min_{\sigma_2 \in \Sigma_2} \max_{s_1 \in S_1} U(s_1, \sigma_2)
\]

\begin{defbox}{Optimal Strategies}
  \begin{itemize}
      \item For Player 1, a \textbf{max-min strategy} is a mixed strategy $\sigma_1^*$ such that:
      \[
        U(\sigma_1^*, \sigma_2) \geq V \quad \forall \sigma_2 \in \Sigma_2
      \]
      \item For Player 2, a \textbf{min-max strategy} is a mixed strategy $\sigma_2^*$ such that:
      \[
        U(\sigma_1, \sigma_2^*) \leq V \quad \forall \sigma_1 \in \Sigma_1
      \]
  \end{itemize}
\end{defbox}

% ============================================================
\section{The Graphical Method ($2 \times 2$ Games)}

When Player 1 has exactly two pure strategies, we can find the max-min strategy by graphing the expected payoffs.

\begin{exbox}{A $2 \times 2$ Game Example}
  \begin{center}
  \renewcommand{\arraystretch}{1.3}
  \begin{tabular}{c|cc}
    $P_1 \setminus P_2$ & L ($y$) & R ($1-y$) \\
    \hline
    T ($x$) & 5 & 0 \\
    B ($1-x$) & 3 & 4 \\
  \end{tabular}
  \end{center}
\end{exbox}

Let Player 1 play T with probability $x$ and B with probability $1-x$.
Player 1's expected payoff against P2's pure strategies:
\begin{align*}
  U(x, L) &= 5x + 3(1-x) = 2x + 3 \\
  U(x, R) &= 0x + 4(1-x) = 4 - 4x
\end{align*}

We graph these two lines for $x \in [0,1]$. For any chosen $x$, Player 2 will respond by choosing the column that yields the \textbf{lower} payoff for Player 1. Thus, Player 1's guaranteed payoff for a given $x$ is the \textbf{lower envelope} of the two lines:
\[
  \min \{2x + 3, 4 - 4x\}
\]
Player 1 wants to maximize this minimum (find the peak of the lower envelope). This peak occurs where the two lines intersect:
\[
  2x + 3 = 4 - 4x \implies 6x = 1 \implies x^* = \frac{1}{6}
\]
The unique max-min strategy for Player 1 is $x^* = 1/6$.
The \textbf{Value} of the game is $V = 2(1/6) + 3 = 10/3 = 3\frac{1}{3}$.

\vspace{0.3cm}
By symmetry, to find Player 2's min-max strategy $y$ (probability of L), we graph P1's payoffs against P1's pure strategies:
\begin{align*}
  U(T, y) &= 5y + 0(1-y) = 5y \\
  U(B, y) &= 3y + 4(1-y) = 4 - y
\end{align*}
Player 2 wants to minimize the \textbf{upper envelope} (because P1 will maximize).
The intersection is $5y = 4 - y \implies 6y = 4 \implies y^* = 2/3$.

% ============================================================
\section{The Graphical Method ($2 \times n$ Games)}

The graphical method also works if one player has 2 strategies and the other has $n$.

\begin{exbox}{A $2 \times 3$ Game Example}
  \begin{center}
  \renewcommand{\arraystretch}{1.3}
  \begin{tabular}{c|ccc}
    $P_1 \setminus P_2$ & L & M & R \\
    \hline
    T ($x$) & 2 & 5 & -1 \\
    B ($1-x$) & 1 & -2 & 5 \\
  \end{tabular}
  \end{center}
\end{exbox}

Player 1's expected payoff against P2's three pure strategies:
\begin{align*}
  U(x, L) &= 2x + 1(1-x) = x + 1 \\
  U(x, M) &= 5x - 2(1-x) = 7x - 2 \\
  U(x, R) &= -x + 5(1-x) = 5 - 6x
\end{align*}

Graphing these three lines for $x \in [0,1]$, we look at the \textbf{lower envelope}.
The highest point on the lower envelope occurs at the intersection of the lines for $L$ and $R$. (The line for $M$ is strictly above this intersection point, so it does not restrict Player 1's maximum guaranteed payoff).
\[
  x + 1 = 5 - 6x \implies 7x = 4 \implies x^* = \frac{4}{7}
\]
The unique max-min strategy for Player 1 is $x^* = 4/7$.
The \textbf{Value} of the game is $V = U(x^*, R) = 5 - 6(4/7) = 35/7 - 24/7 = 11/7$.

\subsection{Finding Player 2's Mixed Strategy}
Player 2 must construct a mixed strategy $\sigma_2^* = (\sigma_2(L), \sigma_2(M), \sigma_2(R))$ such that Player 1 can never get more than $V = 11/7$.
Because the line for pure strategy $M$ is strictly above the value $11/7$ at $x^*$, Player 2 should never play $M$ (it would give P1 a higher payoff). Thus, $\sigma_2(M) = 0$.

Let $\sigma_2(L) = q$ and $\sigma_2(R) = 1-q$.
Player 1's expected payoff against this mix must be $\leq 11/7$ for any pure strategy P1 plays:
\begin{align*}
  \text{If P1 plays T:} \quad & 2q + (-1)(1-q) \leq \frac{11}{7} \implies 3q - 1 \leq \frac{11}{7} \implies 3q \leq \frac{18}{7} \implies q \leq \frac{6}{7} \\
  \text{If P1 plays B:} \quad & 1q + 5(1-q) \leq \frac{11}{7} \implies 5 - 4q \leq \frac{11}{7} \implies 4q \geq \frac{24}{7} \implies q \geq \frac{6}{7}
\end{align*}
Therefore, $q = 6/7$. Player 2's unique min-max strategy is $[6/7 \text{ (L)}, 0 \text{ (M)}, 1/7 \text{ (R)}]$.

\end{document}